\documentclass{article}
\usepackage[english]{babel}
\usepackage[letterpaper,top=2cm,bottom=2cm,left=3cm,right=3cm,marginparwidth=1.75cm]{geometry}
\usepackage{enumitem}
\usepackage{amsmath}
\usepackage[colorlinks=true, allcolors=blue]{hyperref}
\usepackage{minted}
\usepackage{float}
\usepackage{graphicx}
\usepackage{multirow}
\usepackage{subcaption}
\usepackage{booktabs}
\usepackage[utf8]{inputenc}
\usepackage{textgreek}
\usepackage[toc,page]{appendix}
\usepackage{xcolor}
\usepackage{nameref}
\usepackage{url}
\usepackage[T1]{fontenc}
\usepackage{todonotes}
\title{Final Project Report}
\author{Aaron Bussche \and  David Magaram\\[0.3em] Group 28 \and Jason Kasdiran}
\begin{document}
	\maketitle

\section*{Introduction}
Final Report: Aim for a max 10 pages document, excluding references and contributions of team members. We will not subtract points for small deviations from the 10 pages, but keep in mind that if your report is much longer, there is likely an issue with conciseness / data presentation somewhere!


\section*{Problem Statement}
Problem statement outlining the specific goals of the work. Use the project proposal assignment and feedback as input.




Neural Cellular Automata (NCAs) are small, shared-weight networks that grow images from a single seed pixel and can regenerate after damage, making them a promising model for robust, self-organizing systems. Their perceived robustness has led to NCAs being deployed as a defensive component in larger architectures: Xu et al. \cite{xu2024adanca} use NCA layers as plug-and-play adaptors inside Vision Transformers to make them more resistant against adversarial samples. Whether NCAs themselves can be subverted is therefore not just of theoretical interest but relevant to the security of systems that rely on them.

Randazzo et al. \cite{randazzo2021adversarial} and Cavuoti et al. \cite{cavuoti2022adversarial} demonstrated targeted attacks that, for example, can make a generated lizard image lose its tail or change color. However, every published attack on NCAs relies on white-box, gradient-based methods that require full access to the model's weights and backpropagation through the rollout. Robustness claims are so far tested only against this strong threat model. Additionally, focusing on the nature-inspired roots of NCAs, a white-box is not what a realistic adversary, for instance in a biological or bioengineering setting, would plausibly have access to.

We therefore ask: \textit{can a query-only, black-box attacker using evolutionary search reproduce the same targeted reprogramming of a Growing NCA, and how close does it get to the gradient-based ceiling?} We attack the NCA via a $16 \times 16$ state-perturbation matrix applied to every cell's 16-dimensional hidden state at every timestep, the same attack surface used by Randazzo et al. Our attacker observes only the final RGB output and optimizes the L2 distance to a chosen target image using CMA-ES, without ever accessing model weights, hidden channels, or gradient information.

If Growing NCAs can be influenced by this attack in a similar way as through white-box approaches, they may be less robust than priorly assumed. Otherwise, they may be safe from more realistic attack types and to be used in public-facing, but not open-weights, ML systems.  




\section*{Related Work}
Review of related work and how your project fits in with this. Use the project proposal assignment and feedback as input.



Mordvintsev et al. \cite{mordvintsev2020growing} introduced the Growing NCA model and showed that, through a sample pool training strategy and exposure to damage during training, the learned local update rules become robust enough to regenerate damaged parts. Due to this robustness, NCAs have also been used defensively: Xu et al. \cite{xu2024adanca} inserted NCA layers between vision transformer blocks to serve as an armor against adversarial attacks.

On the offensive side, Randazzo et al. \cite{randazzo2021adversarial} investigated white-box vulnerabilities of Growing NCAs and showed that, while they are more resistant to local injections, they can still be reprogrammed using global state perturbations applied to every cell at every step. Similarly, Cavuoti et al. \cite{cavuoti2022adversarial} demonstrated that injecting a small number of adversarial cells into the grid can drastically alter the grown image. Both studies rely on gradient-based optimization with full access to model weights.

Covariance Matrix Adaptation Evolution Strategy (CMA-ES), introduced by Hansen \cite{hansen2001cmaes}, is an algorithm for non-linear, non-convex black-box optimization that maintains a population of candidates and iteratively shifts them toward high-fitness regions. While evolutionary black-box attacks are commonly used to evaluate traditional image classifiers, they have not yet been applied as an adversarial tool against cellular automata.






\section*{Motivation}
Motivation and description of your approach, including: methods, algorithms, evaluation/testing procedures, experimental setup, etc. Ensure that the work is reproducible given this description alone.

\section*{Results}
Results. Ensure to perform a correct experimental analysis (e.g. repeated experiments, statistical analysis of results). Verbally describe your results in this section and make sure to reference the corresponding figures/tables.


\section*{Discussion}
Discussion. A thorough discussion based on the analysis of the results, including limitations and future directions (What additional enhancements could be made? What future work should be done?)


\section*{Conclusion}
Conclusions. Is your line of work worth pursuing?

\newpage
\appendix
\section*{Appendix}


\subsection*{Bibliography}




\subsection*{Author Contribution}
Author contribution: describe each team member contributions to the project.


\section*{Code}
Include a link to a github repository with implementation of algorithms developed and used during the project. The repository should include
source code,
all data necessary to run the program (think instructions for installing external software, package versions/virtual environment, etc);
a short manual describing how to use/run the program;
a sample run, screenshot, or other indication of system behavior.



\newpage
\begin{thebibliography}{9}

\bibitem{mordvintsev2020growing}
A.~Mordvintsev, E.~Randazzo, E.~Niklasson, and M.~Levin,
``Growing neural cellular automata,''
\textit{Distill}, vol.~5, no.~2, e23, 2020.
[Online]. Available: \url{https://distill.pub/2020/growing-ca}
doi: \texttt{10.23915/distill.00023}.


\bibitem{xu2024adanca}
Y.~Xu, T.~Zhang, and S.~Süsstrunk,
``AdaNCA: Neural cellular automata as adaptors for more robust vision transformer,''
\textit{Advances in Neural Information Processing Systems}, vol.~37, pp.~25709--25746, 2024.


\bibitem{randazzo2021adversarial}
E.~Randazzo, A.~Mordvintsev, E.~Niklasson, and M.~Levin,
``Adversarial reprogramming of neural cellular automata,''
\textit{Distill}, vol.~6, no.~5, e00027-004, 2021.
[Online]. Available: \url{https://distill.pub/selforg/2021/adversarial}
doi: \texttt{10.23915/distill.00027.004}.


\bibitem{cavuoti2022adversarial}
L.~Cavuoti, F.~Sacco, E.~Randazzo, and M.~Levin,
``Adversarial takeover of neural cellular automata,''
in \textit{Artificial Life Conference Proceedings 34}, 2022, p.~38.
MIT Press.


\bibitem{hansen2001cmaes}
N.~Hansen and A.~Ostermeier,
``Completely derandomized self-adaptation in evolution strategies,''
\textit{Evolutionary Computation}, vol.~9, no.~2, pp.~159--195, 2001.


\bibitem{greydanus2022growing}
S.~Greydanus,
``Studying growth with neural cellular automata,''
blog post and PyTorch implementation, 2022.
[Online]. Available: \url{https://greydanus.github.io/2022/05/24/studying-growth}


\end{thebibliography}


\end{document}